\documentclass[a4paper,11pt]{amsart}

\usepackage[margin=1in]{geometry}
\usepackage{amsmath,amssymb,mathtools}
\usepackage[T1]{fontenc}
\usepackage{microtype}

\newcommand{\Z}{\mathbb Z}
\newcommand{\Q}{\mathbb Q}
\newcommand{\R}{\mathbb R}
\newcommand{\C}{\mathbb C}
\newcommand{\Pj}{\mathbb P}
\newcommand{\Comp}{\operatorname{Comp}}
\newcommand{\Sp}{\operatorname{Sp}}
\newcommand{\Sym}{\operatorname{Sym}}
\newcommand{\Lie}{\operatorname{Lie}}
\newcommand{\GL}{\operatorname{GL}}

\newtheorem{theorem}{Theorem}[section]
\newtheorem{proposition}[theorem]{Proposition}
\newtheorem{criterion}[theorem]{Criterion}
\newtheorem{lemma}[theorem]{Lemma}
\newtheorem{corollary}[theorem]{Corollary}
\theoremstyle{remark}
\newtheorem{remark}[theorem]{Remark}

\numberwithin{equation}{section}

\title{A Conditional Positivity Reduction}
\date{}

\begin{document}

\begin{abstract}
We isolate an elementary determinant positivity criterion which implies thinness for the hypergeometric tower generated by the companion matrices of \((x-1)^{2n-2}\Phi_6(x)\) and \(x^{2n}+1\). The paper does not prove this positivity criterion. Instead, it proves a self-contained reduction from the criterion to projective freeness, Zariski density, and infinite index in the ambient integral symplectic group. The group is put into a uniform signed-cyclic transvection normal form, and every cyclically reduced mixed word is detected by an upper-Hessenberg determinant with an anti-periodic kernel. Thus a proof of the folded-Hessenberg positivity criterion would imply that the projectivized group is \(\Z*(\Z/2n\Z)\) for every \(n\geq3\).
\end{abstract}

\maketitle

\section{The positivity criterion}

Let \(d\geq6\) be even. Define
\[
Q_d(z)
=
(1+z)^{d-2}(1+z+z^2)
=
\sum_{r=0}^{d}q_rz^r.
\]
Equivalently,
\[
q_r=
\binom{d-2}{r}
+\binom{d-2}{r-1}
+\binom{d-2}{r-2},
\]
where out-of-range binomial coefficients are zero. Hence
\[
q_0=q_d=1,
\qquad
q_r=q_{d-r}>0
\quad(1\leq r\leq d-1).
\]

Define the \(d\)-anti-periodic extension
\[
a_d:\Z\longrightarrow\Z
\]
by
\begin{equation}\label{eq:adef}
a_d(md)=0,
\qquad
a_d(md+r)=(-1)^mq_r
\quad(1\leq r\leq d-1).
\end{equation}
Thus
\[
a_d(r+d)=-a_d(r),
\qquad
a_d(-r)=-a_d(r).
\]

Let
\[
k=(k_1,\ldots,k_\ell)\in\{1,\ldots,d-1\}^{\ell}.
\]
For
\[
\eta=(\eta_1,\ldots,\eta_{\ell-1}),
\]
define the upper-Hessenberg matrix \(H_d(k,\eta)\) by
\begin{equation}\label{eq:Hdef}
[H_d(k,\eta)]_{ij}
=
\begin{cases}
a_d(k_i+\cdots+k_j),&i\leq j,\\[2mm]
\eta_j,&i=j+1,\\[2mm]
0,&i\geq j+2.
\end{cases}
\end{equation}
For \(\ell=1\), this means
\[
H_d((k_1),())=(q_{k_1}).
\]

\begin{criterion}[Folded-Hessenberg positivity]\label{crit:folded-hessenberg}
For every even \(d\geq6\), every \(\ell\geq1\), every
\[
k\in\{1,\ldots,d-1\}^{\ell},
\]
and every
\[
\eta\in[-1,1]^{\ell-1},
\]
one has
\begin{equation}\label{eq:criterion}
\det H_d(k,\eta)>0.
\end{equation}
\end{criterion}

\begin{proposition}\label{prop:sign-version}
Criterion~\ref{crit:folded-hessenberg} is equivalent to its restriction to the vertices
\[
\eta\in\{\pm1\}^{\ell-1}.
\]
\end{proposition}

\begin{proof}
Each parameter \(\eta_j\) appears in exactly one matrix entry, namely in position
\[
(j+1,j).
\]
Therefore \(\det H_d(k,\eta)\) is affine in each variable \(\eta_j\). For \(\eta\in[-1,1]^{\ell-1}\), multiaffine interpolation gives
\[
\det H_d(k,\eta)
=
\sum_{\varepsilon\in\{\pm1\}^{\ell-1}}
\left(
\prod_{j=1}^{\ell-1}
\frac{1+\varepsilon_j\eta_j}{2}
\right)
\det H_d(k,\varepsilon).
\]
The coefficients are nonnegative and sum to one. Thus positivity at all vertices implies positivity throughout the cube. The converse is immediate.
\end{proof}

\begin{remark}
The reduction below requires only the weaker discrete nonvanishing statement
\[
\det H_d\left(k,-\frac1{m_2},\ldots,-\frac1{m_\ell}\right)\neq0
\]
for nonzero integers \(m_2,\ldots,m_\ell\). Criterion~\ref{crit:folded-hessenberg} is a stronger and cleaner formulation.
\end{remark}

\section{Normal form}

Let
\[
f_d(x)=(x-1)^{d-2}\Phi_6(x)=(x-1)^{d-2}(x^2-x+1),
\qquad
g_d(x)=x^d+1.
\]
Let
\[
A_0=\Comp(f_d),
\qquad
B_0=\Comp(g_d),
\qquad
\Gamma_d^{\mathrm{orig}}=\langle A_0,B_0\rangle.
\]
We use the low-degree-first companion convention: if
\[
p(x)=x^d+p_{d-1}x^{d-1}+\cdots+p_1x+p_0,
\]
then
\[
\Comp(p)e_j=e_{j+1}
\quad(0\leq j<d-1),
\qquad
\Comp(p)e_{d-1}
=
-\sum_{j=0}^{d-1}p_je_j.
\]

Since
\[
Q_d(z)=f_d(-z),
\]
the coefficients of \(f_d\) are alternating:
\[
f_d(x)=\sum_{r=0}^{d}(-1)^rq_rx^r.
\]

\begin{lemma}[Signed-cyclic normal form]\label{lem:normal-form}
There is a matrix
\[
P\in M_d(\Z)\cap \GL_d(\Q)
\]
such that
\[
P^{-1}\Gamma_d^{\mathrm{orig}}P=\langle S,T\rangle,
\]
where \(S,T\in\GL_d(\Z)\) are given by
\begin{equation}\label{eq:Saction}
Se_j=
\begin{cases}
-e_{j+1},&0\leq j<d-1,\\
e_0,&j=d-1,
\end{cases}
\end{equation}
and
\begin{equation}\label{eq:Taction}
T(x_0,\ldots,x_{d-1})
=
\left(
x_0+\sum_{r=1}^{d-1}q_rx_r,\,
x_1,\ldots,x_{d-1}
\right).
\end{equation}
In particular,
\[
S^d=-I.
\]
The conjugacy is rational; it is not asserted to be unimodular.
\end{lemma}

\begin{proof}
Let
\[
T_0=B_0A_0^{-1}.
\]
Since
\[
\Gamma_d^{\mathrm{orig}}=\langle A_0,B_0\rangle
=\langle T_0,B_0\rangle,
\]
it suffices to put \(B_0\) and \(T_0\) in the claimed form.

Comparing the last columns of \(A_0\) and \(B_0\) gives
\[
B_0-A_0=-ve_{d-1}^{t},
\]
where
\[
v=\sum_{r=1}^{d-1}(-1)^{r+1}q_re_r.
\]
The constant coefficient of \(f_d\) is one, so the last row of \(A_0^{-1}\) is
\[
-e_0^t.
\]
Indeed, if \(y=A_0x\), then the zeroth coordinate is
\[
y_0=-x_{d-1},
\]
because the constant coefficient of \(f_d\) is one. Hence
\[
x_{d-1}=-y_0.
\]
Therefore
\begin{equation}\label{eq:rankone-T0}
T_0-I=(B_0-A_0)A_0^{-1}=ve_0^t.
\end{equation}

Identify \(\Q^d\) with
\[
\Q[x]/(x^d+1),
\]
where \(B_0\) acts by multiplication by \(x\). Put \(z=-x\). The vector \(v\) is represented modulo \(z^d+1\) by
\[
-\sum_{r=1}^{d-1}q_rz^r
\equiv
-Q_d(z).
\]
The roots of \(Q_d(z)\) are \(-1\) and the two primitive third roots of unity. None of them is a root of \(z^d+1\), since \(d\) is even. Hence
\[
\gcd(Q_d(z),z^d+1)=1.
\]

Define
\[
P=
\bigl[
v,\;(-B_0)v,\;\ldots,\;(-B_0)^{d-1}v
\bigr].
\]
The coprimality just proved implies that these columns form a \(\Q\)-basis. Hence
\[
P\in M_d(\Z)\cap\GL_d(\Q).
\]
The matrix is integral because \(B_0\) and \(v\) are integral, but it need not be unimodular.

In this basis, \(B_0\) acts as \(S\), giving \eqref{eq:Saction}. It remains to compute \(T_0\). By \eqref{eq:rankone-T0}, the coefficient of \(v\) in
\[
(T_0-I)(-B_0)^rv
\]
is the first coordinate of \((-B_0)^rv\). In the polynomial model this is the constant coefficient of
\[
-z^rQ_d(z)
\quad\text{modulo } z^d+1.
\]

For \(r=0\), one must reduce the endpoint terms:
\[
-Q_d(z)
\equiv
-Q_d(z)+(1+z^d)
=
-\sum_{s=1}^{d-1}q_sz^s
\pmod{z^d+1}.
\]
Thus the constant coefficient is zero. For \(1\leq r\leq d-1\), the only term contributing to the constant coefficient after reduction modulo \(z^d+1\) is the term with exponent \(d\), namely
\[
-q_{d-r}z^d\equiv q_{d-r}.
\]
Hence the constant coefficient is
\[
q_{d-r}=q_r.
\]
Therefore \(T_0\) acts as \eqref{eq:Taction}.
\end{proof}

From now until the final lattice-transfer paragraph, write
\[
\Gamma_d=\langle S,T\rangle.
\]

Define an alternating form \(\Omega_d\) by
\begin{equation}\label{eq:Omega}
\Omega_d(e_i,e_j)=
\begin{cases}
q_{j-i},&i<j,\\
-q_{i-j},&i>j,\\
0,&i=j.
\end{cases}
\end{equation}

\begin{proposition}\label{prop:symplectic-form}
The form \(\Omega_d\) is nondegenerate, and both \(S\) and \(T\) preserve it. Moreover,
\begin{equation}\label{eq:T-transvection}
Tx=x+\Omega_d(e_0,x)e_0.
\end{equation}
\end{proposition}

\begin{proof}
The equality
\[
S^t\Omega_dS=\Omega_d
\]
follows directly from \eqref{eq:Saction} and the symmetry
\[
q_r=q_{d-r}.
\]
Formula \eqref{eq:T-transvection} is exactly \eqref{eq:Taction}. Hence \(T\) is a symplectic transvection and preserves \(\Omega_d\).

It remains to prove nondegeneracy. The matrix \(S\) is a signed permutation matrix, so
\[
S^{-t}=S.
\]
Thus
\[
S^t\Omega_dS=\Omega_d
\]
implies
\[
\Omega_dS=S^{-t}\Omega_d=S\Omega_d.
\]
Consequently \(\Omega_d\), regarded as a linear map using the standard coordinates, commutes with \(S\).

The eigenvalues of \(S\) are the distinct roots of
\[
\lambda^d=-1.
\]
For such a root \(\lambda\), define
\[
w_\lambda=
\left(
1,-\lambda^{-1},\lambda^{-2},\ldots,
(-1)^{d-1}\lambda^{-(d-1)}
\right)^t.
\]
Then
\[
Sw_\lambda=\lambda w_\lambda.
\]
Since \(\Omega_dS=S\Omega_d\), the one-dimensional eigenspace spanned by \(w_\lambda\) is \(\Omega_d\)-stable. The eigenvalue of \(\Omega_d\) on \(w_\lambda\) is the zeroth coordinate of \(\Omega_dw_\lambda\), namely
\[
\sum_{r=1}^{d-1}q_r(-\lambda^{-1})^r.
\]
Put
\[
t=-\lambda^{-1}.
\]
Then \(t^d=-1\), and hence
\[
\sum_{r=1}^{d-1}q_rt^r
=
Q_d(t)-q_0-q_dt^d
=
Q_d(t).
\]
This is nonzero, because the roots of \(Q_d\) are \(-1\) and the primitive third roots of unity, none of which satisfies \(t^d=-1\). Thus every eigenvalue of \(\Omega_d\) is nonzero.
\end{proof}

\section{Hessenberg expansion and reduced words}

For \(j\in\Z\), define
\begin{equation}\label{eq:u_j}
u_j=(-1)^jS^je_0.
\end{equation}
Then
\[
u_{j+d}=-u_j.
\]

\begin{lemma}\label{lem:pairing}
For all \(p,q\in\Z\),
\begin{equation}\label{eq:pairing}
\Omega_d(u_p,u_q)=a_d(q-p).
\end{equation}
\end{lemma}

\begin{proof}
First,
\[
\Omega_d(u_p,u_q)
=
(-1)^{p+q}\Omega_d(S^pe_0,S^qe_0).
\]
By \(S\)-invariance of \(\Omega_d\),
\[
\Omega_d(S^pe_0,S^qe_0)
=
\Omega_d(e_0,S^{q-p}e_0).
\]
Since
\[
(-1)^{p+q}=(-1)^{q-p},
\]
we get
\[
\Omega_d(u_p,u_q)
=
\Omega_d(e_0,(-1)^{q-p}S^{q-p}e_0)
=
\Omega_d(u_0,u_{q-p}).
\]
It remains to compute the case \(p=0\).

If
\[
q=md+r,
\qquad
1\leq r\leq d-1,
\]
then
\[
u_q=(-1)^me_r,
\]
and therefore
\[
\Omega_d(u_0,u_q)=(-1)^mq_r=a_d(q).
\]
If \(d\mid q\), then \(u_q=\pm u_0\), and both sides are zero.
\end{proof}

For a vector \(v\), write
\[
\tau_v(x)=x+\Omega_d(v,x)v.
\]
Then
\[
T=\tau_{u_0},
\qquad
\tau_{-v}=\tau_v.
\]

We next record the elementary Hessenberg expansion used in the reduced-word calculation.

\begin{lemma}[Interval expansion]\label{lem:interval-expansion}
Let \(h_{ij}\) be scalars for \(1\leq i\leq j\leq\ell\). Let \(M\) be the upper-Hessenberg matrix
\[
M_{ij}
=
\begin{cases}
h_{ij},&i\leq j,\\
\eta_j,&i=j+1,\\
0,&i\geq j+2.
\end{cases}
\]
Then
\begin{equation}\label{eq:interval-expansion}
\det M
=
\sum_{0=i_0<i_1<\cdots<i_r=\ell}
(-1)^{\ell-r}
\left(
\prod_{j\in\{1,\ldots,\ell-1\}\setminus\{i_1,\ldots,i_{r-1}\}}\eta_j
\right)
\prod_{a=0}^{r-1}
h_{i_a+1,i_{a+1}}.
\end{equation}
\end{lemma}

\begin{proof}
Expand the determinant as a sum over permutations \(\sigma\in S_\ell\). A term is nonzero only if
\[
\sigma(i)\geq i-1
\qquad(1\leq i\leq\ell).
\]
Such a permutation decomposes uniquely into interval cycles. Indeed, if \(b\) is the smallest row not already used, then necessarily
\[
\sigma(b)\geq b.
\]
If \(\sigma(b)=c\), then all rows
\[
b+1,\ldots,c
\]
must map to their immediate predecessors:
\[
j\longmapsto j-1.
\]
Otherwise one of the columns
\[
b,\ldots,c-1
\]
would be missed or used twice. Hence the cycle beginning at \(b\) is exactly the interval cycle
\[
b\longmapsto c,
\qquad
j\longmapsto j-1
\quad(b+1\leq j\leq c).
\]
Repeating with the next unused row gives the interval decomposition uniquely.

Equivalently, choose a chain
\[
0=i_0<i_1<\cdots<i_r=\ell.
\]
For the block
\[
\{i_a+1,\ldots,i_{a+1}\},
\]
the corresponding cycle sends
\[
i_a+1\longmapsto i_{a+1},
\qquad
j\longmapsto j-1
\quad(i_a+2\leq j\leq i_{a+1}).
\]
A block of length one is the fixed point
\[
i_a+1=i_{a+1}.
\]
The block contributes the upper entry
\[
h_{i_a+1,i_{a+1}},
\]
the subdiagonal entries
\[
\eta_{i_a+1},\eta_{i_a+2},\ldots,\eta_{i_{a+1}-1},
\]
and has sign
\[
(-1)^{i_{a+1}-i_a-1}.
\]
Multiplying over all blocks gives total sign
\[
(-1)^{\ell-r}.
\]
The subdiagonal factors used are precisely those \(\eta_j\) for which \(j\) is not one of the interior endpoints
\[
i_1,\ldots,i_{r-1}.
\]
This proves \eqref{eq:interval-expansion}.
\end{proof}

We also need the matching chain expansion for products of transvections.

\begin{lemma}[Transvection-chain expansion]\label{lem:chain-expansion}
Let
\[
m_2,\ldots,m_\ell\in\Z\setminus\{0\}
\]
and let
\[
p_0,\ldots,p_\ell
\]
be integers. Then
\begin{align}
&\Omega_d\left(
u_{p_0},
\tau_{u_{p_1}}^{m_2}
\tau_{u_{p_2}}^{m_3}
\cdots
\tau_{u_{p_{\ell-1}}}^{m_\ell}
u_{p_\ell}
\right)\notag\\
&\quad=
\sum_{0=i_0<i_1<\cdots<i_r=\ell}
\left(
\prod_{a=1}^{r-1}m_{i_a+1}
\right)
\prod_{a=0}^{r-1}
a_d(p_{i_{a+1}}-p_{i_a}).
\label{eq:chain-expansion}
\end{align}
For \(\ell=1\), the product of transvections and the product over \(a=1,\ldots,r-1\) are empty.
\end{lemma}

\begin{proof}
Write
\[
\tau_{u_{p_j}}^{m_{j+1}}
=
I+m_{j+1}N_j,
\qquad
N_j(x)=\Omega_d(u_{p_j},x)u_{p_j}
\quad(1\leq j\leq\ell-1).
\]
Expand the product
\[
(I+m_2N_1)(I+m_3N_2)\cdots(I+m_\ell N_{\ell-1}).
\]
Choose a subset
\[
J=\{i_1<\cdots<i_{r-1}\}
\subseteq\{1,\ldots,\ell-1\},
\]
and set
\[
i_0=0,
\qquad
i_r=\ell.
\]
The corresponding term is
\[
\left(\prod_{a=1}^{r-1}m_{i_a+1}\right)
\Omega_d\left(
u_{p_0},
N_{i_1}N_{i_2}\cdots N_{i_{r-1}}u_{p_\ell}
\right).
\]
Since
\[
N_j(x)=\Omega_d(u_{p_j},x)u_{p_j},
\]
this pairing equals
\[
\left(\prod_{a=1}^{r-1}m_{i_a+1}\right)
\Omega_d(u_{p_0},u_{p_{i_1}})
\Omega_d(u_{p_{i_1}},u_{p_{i_2}})
\cdots
\Omega_d(u_{p_{i_{r-1}}},u_{p_\ell}).
\]
By Lemma~\ref{lem:pairing}, this is exactly the summand in \eqref{eq:chain-expansion}. Summing over all subsets \(J\) proves the formula.
\end{proof}

Let
\[
w=t^{m_1}s^{k_1}\cdots t^{m_\ell}s^{k_\ell}
\]
be a cyclically reduced mixed word in
\[
\langle t\rangle *\langle s\mid s^d=1\rangle,
\]
so
\[
m_i\in\Z\setminus\{0\},
\qquad
1\leq k_i\leq d-1.
\]
Put
\[
p_0=0,
\qquad
p_j=k_1+\cdots+k_j.
\]

Under the homomorphism
\[
t\mapsto \overline T,
\qquad
s\mapsto \overline S,
\]
a linear lift of the image of \(w\) is
\begin{equation}\label{eq:word-lift}
W=
\tau_{u_{p_0}}^{m_1}
\tau_{u_{p_1}}^{m_2}
\cdots
\tau_{u_{p_{\ell-1}}}^{m_\ell}
S^{p_\ell}.
\end{equation}
Indeed, for every integer \(p\),
\[
S^p\tau_{u_0}^mS^{-p}
=
\tau_{S^pu_0}^m
=
\tau_{u_p}^m,
\]
because
\[
S^pu_0=\pm u_p
\]
and
\[
\tau_{-v}=\tau_v.
\]
Pushing all powers of \(S\) to the right in
\[
T^{m_1}S^{k_1}\cdots T^{m_\ell}S^{k_\ell}
\]
therefore gives \eqref{eq:word-lift}.

\begin{lemma}[Reduced-word determinant]\label{lem:word-determinant}
Let
\[
\eta_j=-\frac1{m_{j+1}}
\qquad(1\leq j\leq \ell-1).
\]
Then
\begin{equation}\label{eq:word-determinant}
\Omega_d(u_0,We_0)
=
(-1)^{p_\ell}
\left(\prod_{i=2}^{\ell}m_i\right)
\det H_d(k,\eta).
\end{equation}
For \(\ell=1\), the empty product is interpreted as one and
\[
H_d(k,\eta)=(q_{k_1}).
\]
\end{lemma}

\begin{proof}
Since
\[
S^{p_\ell}e_0=(-1)^{p_\ell}u_{p_\ell},
\]
and since \(\tau_{u_0}\) fixes the pairing with \(u_0\), we have
\[
\Omega_d(u_0,We_0)
=
(-1)^{p_\ell}
\Omega_d\left(
u_0,\,
\tau_{u_{p_1}}^{m_2}
\cdots
\tau_{u_{p_{\ell-1}}}^{m_\ell}
u_{p_\ell}
\right).
\]
By Lemma~\ref{lem:chain-expansion}, the latter pairing is
\begin{equation}\label{eq:chain-side}
\sum_{0=i_0<i_1<\cdots<i_r=\ell}
\left(
\prod_{a=1}^{r-1}m_{i_a+1}
\right)
\prod_{a=0}^{r-1}
a_d(p_{i_{a+1}}-p_{i_a}).
\end{equation}

On the determinant side, apply Lemma~\ref{lem:interval-expansion} with
\[
h_{ij}=a_d(p_j-p_{i-1}).
\]
The same chain contributes
\begin{equation}\label{eq:block-side}
(-1)^{\ell-r}
\left(
\prod_{j\in\{1,\ldots,\ell-1\}\setminus\{i_1,\ldots,i_{r-1}\}}\eta_j
\right)
\prod_{a=0}^{r-1}
a_d(p_{i_{a+1}}-p_{i_a}).
\end{equation}
Substituting
\[
\eta_j=-\frac1{m_{j+1}}
\]
and multiplying by
\[
\prod_{i=2}^{\ell}m_i
\]
turns \eqref{eq:block-side} into the corresponding summand of \eqref{eq:chain-side}: the sign
\[
(-1)^{\ell-r}
\]
is cancelled by the \(\ell-r\) minus signs coming from the \(\eta_j\)'s, and the remaining \(m\)-factors are precisely
\[
\prod_{a=1}^{r-1}m_{i_a+1}.
\]
This proves \eqref{eq:word-determinant}.
\end{proof}

\section{Conditional projective freeness}

\begin{theorem}[Criterion implies projective freeness]\label{thm:criterion-free}
Assume Criterion~\ref{crit:folded-hessenberg} for the given even integer \(d\geq6\). Then
\begin{equation}\label{eq:projective-free-product}
\Pj\Gamma_d\cong \Z*(\Z/d\Z).
\end{equation}
\end{theorem}

\begin{proof}
The projective order of \(S\) is exactly \(d\), because
\[
S^d=-I
\]
and no smaller positive power of the signed cyclic shift is scalar. The element \(T\) has infinite projective order, because
\[
T^m=I+m(T-I)
\]
is not scalar for every \(m\neq0\).

Thus there is a natural surjective homomorphism
\[
\Z*(\Z/d\Z)\longrightarrow \Pj\Gamma_d.
\]
It remains to prove injectivity. Suppose a nontrivial reduced word lies in the kernel. If it belongs to one free factor, this contradicts the orders just computed. Otherwise, since the kernel is normal, any cyclic conjugate of the given kernel element is again in the kernel. Thus we may assume the reduced mixed word is cyclically reduced and has the form
\[
w=t^{m_1}s^{k_1}\cdots t^{m_\ell}s^{k_\ell},
\]
with
\[
m_i\neq0,
\qquad
1\leq k_i\leq d-1.
\]
Let \(W\) be the lift \eqref{eq:word-lift}. Since
\[
\eta_j=-\frac1{m_{j+1}}\in[-1,1],
\]
Criterion~\ref{crit:folded-hessenberg} gives
\[
\det H_d(k,\eta)>0.
\]
By Lemma~\ref{lem:word-determinant},
\[
\Omega_d(u_0,We_0)\neq0.
\]
But if \(W\) were scalar, then
\[
\Omega_d(u_0,We_0)
=
\lambda\Omega_d(u_0,e_0)
=
0.
\]
Hence \(W\) is not scalar. This contradiction proves injectivity.
\end{proof}

\section{Zariski density and infinite index}

All Zariski closures in the next proposition are taken in the real algebraic group
\[
\Sp(\Omega_d)_{\R}.
\]
Equivalently, after extension of scalars to \(\C\), the complex Zariski closure is
\[
\Sp(\Omega_d)_{\C}.
\]

\begin{proposition}[Zariski density]\label{prop:zariski-density}
For every even \(d\geq6\),
\[
\overline{\Gamma_d}^{\,Z}=\Sp(\Omega_d).
\]
\end{proposition}

\begin{proof}
For \(0\leq i\leq d-1\), put
\[
v_i=S^ie_0,
\qquad
T_i=S^iTS^{-i}.
\]
Then
\[
T_i=I+N_i,
\qquad
N_i(x)=\Omega_d(v_i,x)v_i.
\]
The vectors \(v_0,\ldots,v_{d-1}\) form a basis. Since
\[
N_i^2=0,
\]
we have
\[
T_i^m=I+mN_i.
\]
The integers are Zariski dense in the affine line, so the Zariski closure of \(\langle T_i\rangle\) contains
\[
I+tN_i
\qquad(t\in\R).
\]
Thus
\[
N_i\in\Lie(\overline{\Gamma_d}^{\,Z}).
\]

For \(i\neq j\), one has
\[
\Omega_d(v_i,v_j)=\pm q_r\neq0
\]
for some \(1\leq r\leq d-1\). Direct calculation gives
\[
[N_i,N_j](x)
=
\Omega_d(v_i,v_j)
\left(
\Omega_d(v_j,x)v_i+\Omega_d(v_i,x)v_j
\right).
\]
For any \(u,v\), define
\[
B_{u,v}(x)=\Omega_d(u,x)v+\Omega_d(v,x)u.
\]
The map
\[
\Sym^2(\R^d)\longrightarrow \mathfrak{sp}(\Omega_d),
\qquad
u\odot v\longmapsto B_{u,v},
\]
is an isomorphism. Indeed, after identifying \(\R^d\) with its dual using the nondegenerate form \(\Omega_d\), this is the usual isomorphism between symmetric bilinear forms and symplectic endomorphisms; equivalently, it is injective and both spaces have dimension \(d(d+1)/2\).

The elements \(N_i\) give the diagonal tensors because
\[
B_{v_i,v_i}=2N_i,
\]
and the commutators \([N_i,N_j]\) give nonzero multiples of \(B_{v_i,v_j}\) for \(i\neq j\). Since the characteristic is zero, the factor \(2\) is harmless. Hence
\[
\Lie(\overline{\Gamma_d}^{\,Z})=\mathfrak{sp}(\Omega_d).
\]
A Zariski-closed subgroup of a smooth algebraic group with the same Lie algebra has the same identity component. Since \(\Sp(\Omega_d)\) is connected, this forces the Zariski closure to be all of \(\Sp(\Omega_d)\).
\end{proof}

\begin{lemma}\label{lem:abelian-free-product}
Every abelian subgroup of
\[
\Z*(\Z/d\Z)
\]
is cyclic.
\end{lemma}

\begin{proof}
Let
\[
G=\Z*(\Z/d\Z).
\]
The group \(G\) acts without inversions on its Bass--Serre tree. Vertex stabilizers are conjugates of the two free factors, and edge stabilizers are trivial.

Let \(A\leq G\) be abelian. If every element of \(A\) is elliptic, then \(A\) fixes a vertex: commuting elliptic automorphisms of a tree have a common fixed point. Hence \(A\) is conjugate into a vertex stabilizer, so \(A\) is a subgroup of either \(\Z\) or \(\Z/d\Z\), and is therefore cyclic.

If \(A\) contains a hyperbolic element \(g\), then every element of \(A\) preserves the axis of \(g\). Since \(A\) is abelian, no element of \(A\) can reverse this axis; a reversal would conjugate the translation \(g\) to \(g^{-1}\). Thus \(A\) acts on the axis by translations. The kernel of this action fixes the axis pointwise, hence fixes an edge, and is trivial because edge stabilizers are trivial. Therefore \(A\) embeds in the infinite cyclic group of translations of the axis. Hence \(A\) is cyclic.
\end{proof}

\begin{proposition}[Infinite index under the criterion]\label{prop:infinite-index}
Assume Criterion~\ref{crit:folded-hessenberg} for the given even integer \(d\geq6\). Then
\[
[\Sp(\Omega_d,\Z):\Gamma_d]=\infty.
\]
\end{proposition}

\begin{proof}
First, \(\Sp(\Omega_d,\Z)\) contains a copy of \(\Z^2\). Let
\[
a=e_0,
\qquad
b=q_2e_1-q_1e_2.
\]
Then \(a\) and \(b\) are independent integral vectors and
\[
\Omega_d(a,b)=q_2q_1-q_1q_2=0.
\]
Therefore the integral transvections
\[
\tau_a(x)=x+\Omega_d(a,x)a,
\qquad
\tau_b(x)=x+\Omega_d(b,x)b
\]
commute. They generate \(\Z^2\): if
\[
\tau_a^r\tau_b^s=I,
\]
then
\[
r\,\Omega_d(a,x)a+s\,\Omega_d(b,x)b=0
\]
for every \(x\), and nondegeneracy of \(\Omega_d\), together with independence of \(a,b\), forces
\[
r=s=0.
\]

Every finite-index subgroup of \(\Sp(\Omega_d,\Z)\) intersects this \(\Z^2\) in a finite-index subgroup, hence contains another copy of \(\Z^2\).

On the other hand, Theorem~\ref{thm:criterion-free} gives
\[
\Pj\Gamma_d\cong \Z*(\Z/d\Z).
\]
By Lemma~\ref{lem:abelian-free-product}, every abelian subgroup of \(\Pj\Gamma_d\) is cyclic. The kernel of
\[
\Gamma_d\longrightarrow \Pj\Gamma_d
\]
is contained in \(\{\pm I\}\), because a scalar symplectic matrix has scalar square equal to one.

Indeed, if \(A\leq\Gamma_d\) were isomorphic to \(\Z^2\), then its image in \(\Pj\Gamma_d\) would be an abelian subgroup, hence cyclic. But
\[
A\cap\{\pm I\}
\]
is finite, so the projection
\[
A\longrightarrow\Pj\Gamma_d
\]
has finite kernel. Therefore the image of \(A\) has rank two, not rank at most one, a contradiction. Thus \(\Gamma_d\) contains no copy of \(\Z^2\), and hence cannot have finite index in \(\Sp(\Omega_d,\Z)\).
\end{proof}

\section{Transfer to the original lattice}

Let \(P\) be the integral rational conjugating matrix from Lemma~\ref{lem:normal-form}. Thus
\[
P^{-1}\Gamma_d^{\mathrm{orig}}P=\Gamma_d.
\]
The lattice
\[
P\Z^d
\]
is a finite-index sublattice of \(\Z^d\), because
\[
P\in M_d(\Z)\cap\GL_d(\Q).
\]

For the original companion-matrix representation, define the rational alternating form
\begin{equation}\label{eq:Omega-orig}
\Omega_d^{\mathrm{orig}}
=
P^{-t}\Omega_dP^{-1}.
\end{equation}
Choose a nonzero integer \(c_d\) such that
\begin{equation}\label{eq:Omega-orig-integral}
\widetilde\Omega_d^{\mathrm{orig}}
=
c_d\Omega_d^{\mathrm{orig}}
\end{equation}
has integral entries. We define the ambient integral symplectic group for the original representation using
\[
\widetilde\Omega_d^{\mathrm{orig}}.
\]
This definition is independent of the choice of \(c_d\), since multiplying the form by a nonzero scalar does not change the equation
\[
M^t\Omega M=\Omega.
\]
Thus
\[
\Sp(\widetilde\Omega_d^{\mathrm{orig}},\Z)
=
\{M\in\GL_d(\Z):M^t\widetilde\Omega_d^{\mathrm{orig}}M
=\widetilde\Omega_d^{\mathrm{orig}}\}.
\]

The group
\[
P\Sp(\Omega_d,\Z)P^{-1}
\]
is the group of \(\Omega_d^{\mathrm{orig}}\)-preserving automorphisms of the commensurable lattice \(P\Z^d\). Since \(P\Z^d\) and \(\Z^d\) are commensurable lattices, the two arithmetic groups
\[
P\Sp(\Omega_d,\Z)P^{-1}
\quad\text{and}\quad
\Sp(\widetilde\Omega_d^{\mathrm{orig}},\Z)
\]
are commensurable. Indeed, after replacing the two lattices by their intersection, the subgroup preserving both lattices has finite index in each corresponding lattice automorphism group.

Therefore Zariski density is unchanged by passing between the conjugated model and the original companion-matrix model. Infinite index transfers as follows. Suppose the original group
\[
\Gamma_d^{\mathrm{orig}}
\]
had finite index in
\[
\Sp(\widetilde\Omega_d^{\mathrm{orig}},\Z).
\]
Since this arithmetic group is commensurable with
\[
P\Sp(\Omega_d,\Z)P^{-1},
\]
the subgroup
\[
P\Gamma_dP^{-1}=\Gamma_d^{\mathrm{orig}}
\]
would have finite index in
\[
P\Sp(\Omega_d,\Z)P^{-1}.
\]
Conjugating by \(P^{-1}\) would imply that \(\Gamma_d\) has finite index in
\[
\Sp(\Omega_d,\Z),
\]
contradicting Proposition~\ref{prop:infinite-index}.

\begin{theorem}[Conditional thinness]\label{thm:conditional-thinness}
Assume Criterion~\ref{crit:folded-hessenberg} for every even \(d\geq6\). Then, for every \(n\geq3\), the group
\[
\left\langle
\Comp\!\left((x-1)^{2n-2}\Phi_6(x)\right),
\Comp(x^{2n}+1)
\right\rangle
\]
is thin in its ambient integral symplectic group. More precisely, its projectivization is
\[
\Z*(\Z/2n\Z).
\]
\end{theorem}

\begin{proof}
Set
\[
d=2n.
\]
Lemma~\ref{lem:normal-form} rationally conjugates the original companion-matrix group to
\[
\Gamma_d=\langle S,T\rangle.
\]
Proposition~\ref{prop:symplectic-form} places \(\Gamma_d\) inside
\[
\Sp(\Omega_d,\Z).
\]

By Theorem~\ref{thm:criterion-free},
\[
\Pj\Gamma_d\cong\Z*(\Z/d\Z).
\]
By Proposition~\ref{prop:zariski-density},
\[
\overline{\Gamma_d}^{\,Z}=\Sp(\Omega_d).
\]
By Proposition~\ref{prop:infinite-index},
\[
[\Sp(\Omega_d,\Z):\Gamma_d]=\infty.
\]
Thus \(\Gamma_d\) is thin in the conjugated model.

The lattice-transfer discussion above identifies the original companion-matrix group with a rational conjugate of \(\Gamma_d\) and identifies its ambient integral symplectic group with a commensurable arithmetic group. Therefore the original companion-matrix group is Zariski dense and infinite-index in its ambient integral symplectic group.
\end{proof}

\begin{remark}
Theorem~\ref{thm:conditional-thinness} is conditional. The paper proves the reduction from folded-Hessenberg positivity to thinness, but it does not prove Criterion~\ref{crit:folded-hessenberg}.
\end{remark}

\end{document}